
\documentclass{article}
\usepackage{colortbl}
\usepackage{makecell}
\usepackage{multirow}
\usepackage{supertabular}

\begin{document}

\newcounter{utterance}

\centering \large Interaction Transcript for game `cladder', experiment `full\_v1.5\_default', episode 1040 with refineonly{-}llama{-}3.1{-}8b{-}instruct.
\vspace{24pt}

{ \footnotesize  \setcounter{utterance}{1}
\setlength{\tabcolsep}{0pt}
\begin{supertabular}{c@{$\;$}|p{.15\linewidth}@{}p{.15\linewidth}p{.15\linewidth}p{.15\linewidth}p{.15\linewidth}p{.15\linewidth}}
    \# & \multicolumn{2}{c}{Player} && \multicolumn{2}{c}{Game Master} \\
    \hline

    \theutterance \stepcounter{utterance}  
    & & & \multicolumn{4}{p{0.6\linewidth}}{
        \cellcolor[rgb]{0.9,0.9,0.9}{
            \makecell[{{p{\linewidth}}}]{
                \texttt{\tiny{[P1$\langle$GM]}}
                \texttt{You are an expert in causal inference. The following question is not a typical commonsense query, but rather a meticulously designed question created by a professor specializing in causal inference, intended to assess the students' mastery of the course content.} \\
\texttt{Imagine a self{-}contained, hypothetical world with only the following conditions, and without any unmentioned factors or causal relationships: Demand has a direct effect on supply and price. Yield per acre has a direct effect on supply. Supply has a direct effect on price. Demand is unobserved. The overall probability of increased supply is 42\%. The probability of reduced supply and increased price is 42\%. The probability of increased supply and increased price is 22\%. Is the chance of increased price smaller when observing increased supply?} \\
\texttt{Start your answer with 'Yes' or 'No', followed by additional reasoning or evidence to support your explanation.} \\
            }
        }
    }
     \\ \\

    \theutterance \stepcounter{utterance}  
    & & \multicolumn{4}{p{0.6\linewidth}}{
        \cellcolor[rgb]{0.95,0.95,0.95}{
            \makecell[{{p{\linewidth}}}]{
                \texttt{\tiny{[P1$\leftrightarrow$P]}}
                \texttt{{[}DRAFT{]}} \\
\texttt{No.} \\
\\ 
\texttt{To answer this question, we need to apply the concept of conditional probability and the causal relationships given in the problem. Let's denote the variables as follows:} \\
\\ 
\texttt{{-} D: Demand} \\
\texttt{{-} S: Supply} \\
\texttt{{-} P: Price} \\
\texttt{{-} Y: Yield per acre} \\
\\ 
\texttt{We are given the following causal relationships:} \\
\\ 
\texttt{{-} D → S (Demand affects supply)} \\
\texttt{{-} Y → S (Yield per acre affects supply)} \\
\texttt{{-} S → P (Supply affects price)} \\
\\ 
\texttt{We are also given the following conditional probabilities:} \\
\\ 
\texttt{{-} P(S ↑) = 0.42 (Probability of increased supply)} \\
\texttt{{-} P(S ↓, P ↑) = 0.42 (Probability of reduced supply and increased price)} \\
\texttt{{-} P(S ↑, P ↑) = 0.22 (Probability of increased supply and increased price)} \\
\\ 
\texttt{We want to find P(P ↑ | S ↑), the probability of increased price given increased supply.} \\
\\ 
\texttt{Using the law of total probability, we can write:} \\
\\ 
\texttt{P(P ↑ | S ↑) = P(P ↑, S ↑ | S ↑) + P(P ↑, S ↓ | S ↑)} \\
\\ 
\texttt{Since S ↑ is given, we can simplify the expression:} \\
\\ 
\texttt{P(P ↑ | S ↑) = P(P ↑, S ↑) / P(S ↑) + P(P ↑, S ↓) / P(S ↑)} \\
\\ 
\texttt{We know that P(S ↑) = P(P ↑, S ↑) + P(P ↑, S ↓) + P(P ↓, S ↑) + P(P ↓, S ↓)} \\
\\ 
\texttt{Substituting this expression into the previous one, we get:} \\
\\ 
\texttt{P(P ↑ | S ↑) = (P(P ↑, S ↑) + P(P ↓, S ↑)) / (P(P ↑, S ↑) + P(P ↑, S ↓) + P(P ↓, S ↑) + P(P ↓, S ↓))} \\
\\ 
\texttt{We are given P(S ↑, P ↑) = 0.22 and P(S ↓, P ↑) = 0.42. We can rewrite these expressions as:} \\
\\ 
\texttt{P(P ↑, S ↑) = 0.22} \\
\texttt{P(P ↓, S ↑) = 0.42 {-} 0.22 = 0.20} \\
\\ 
\texttt{We also know that P(S ↑) = 0.42 and P(S ↓) = 1 {-} P(S ↑) = 0.58.} \\
\\ 
\texttt{Using the law of total probability, we can write:} \\
\\ 
\texttt{P(P ↑) = P(P ↑, S ↑) + P(P ↑, S ↓) = 0.22 + P(P ↑, S ↓)} \\
\\ 
\texttt{We can also write:} \\
\\ 
\texttt{P(P ↑) = P(P ↓, S ↑) + P(P ↓, S ↓) + P(P ↑, S ↑) + P(P ↑, S ↓)} \\
\\ 
\texttt{Substituting the expressions for P(P ↓, S ↑) and P(P ↑, S ↑), we get:} \\
\\ 
\texttt{P(P ↑) = 0.20 + P} \\
            }
        }
    }
    & \\ \\

    \theutterance \stepcounter{utterance}  
    & \multicolumn{4}{p{0.6\linewidth}}{
        \cellcolor[rgb]{0.561,0.737,0.561}{
            \makecell[{{p{\linewidth}}}]{
                \texttt{\tiny{[P1$\rangle$GM]}}
                \texttt{No.} \\
\\ 
\texttt{To answer this question, we need to apply the concept of conditional probability and the causal relationships given in the problem. Let's denote the variables as follows:} \\
\\ 
\texttt{{-} D: Demand} \\
\texttt{{-} S: Supply} \\
\texttt{{-} P: Price} \\
\texttt{{-} Y: Yield per acre} \\
\\ 
\texttt{We are given the following causal relationships:} \\
\\ 
\texttt{{-} D → S (Demand affects supply)} \\
\texttt{{-} Y → S (Yield per acre affects supply)} \\
\texttt{{-} S → P (Supply affects price)} \\
\\ 
\texttt{We are also given the following conditional probabilities:} \\
\\ 
\texttt{{-} P(S ↑) = 0.42 (Probability of increased supply)} \\
\texttt{{-} P(S ↓, P ↑) = 0.42 (Probability of reduced supply and increased price)} \\
\texttt{{-} P(S ↑, P ↑) = 0.22 (Probability of increased supply and increased price)} \\
\\ 
\texttt{We want to find P(P ↑ | S ↑), the probability of increased price given increased supply.} \\
\\ 
\texttt{Using the law of total probability, we can write:} \\
\\ 
\texttt{P(P ↑ | S ↑) = P(P ↑, S ↑ | S ↑) + P(P ↑, S ↓ | S ↑)} \\
\\ 
\texttt{Since S ↑ is given, we can simplify the expression:} \\
\\ 
\texttt{P(P ↑ | S ↑) = P(P ↑, S ↑) / P(S ↑) + P(P ↑, S ↓) / P(S ↑)} \\
\\ 
\texttt{We know that P(S ↑) = P(P ↑, S ↑) + P(P ↑, S ↓) + P(P ↓, S ↑) + P(P ↓, S ↓)} \\
\\ 
\texttt{Substituting this expression into the previous one, we get:} \\
\\ 
\texttt{P(P ↑ | S ↑) = (P(P ↑, S ↑) + P(P ↓, S ↑)) / (P(P ↑, S ↑) + P(P ↑, S ↓) + P(P ↓, S ↑) + P(P ↓, S ↓))} \\
\\ 
\texttt{We are given P(S ↑, P ↑) = 0.22 and P(S ↓, P ↑) = 0.42. We can rewrite these expressions as:} \\
\\ 
\texttt{P(P ↑, S ↑) = 0.22} \\
\texttt{P(P ↓, S ↑) = 0.42 {-} 0.22 = 0.20} \\
\\ 
\texttt{We also know that P(S ↑) = 0.42 and P(S ↓) = 1 {-} P(S ↑) = 0.58.} \\
\\ 
\texttt{Using the law of total probability, we can write:} \\
\\ 
\texttt{P(P ↑) = P(P ↑, S ↑) + P(P ↑, S ↓) = 0.22 + P(P ↑, S ↓)} \\
\\ 
\texttt{We can also write:} \\
\\ 
\texttt{P(P ↑) = P(P ↓, S ↑) + P(P ↓, S ↓) + P(P ↑, S ↑) + P(P ↑, S ↓)} \\
\\ 
\texttt{Substituting the expressions for P(P ↓, S ↑) and P(P ↑, S ↑), we get:} \\
\\ 
\texttt{P(P ↑) = 0.20 + P} \\
            }
        }
    }
    & & \\ \\

    \theutterance \stepcounter{utterance}  
    & & & \multicolumn{2}{p{0.3\linewidth}}{
        \cellcolor[rgb]{0.9,0.9,0.9}{
            \makecell[{{p{\linewidth}}}]{
                \texttt{\tiny{[GM$|$GM]}}
                \texttt{No} \\
            }
        }
    }
    & & \\ \\

    \theutterance \stepcounter{utterance}  
    & & & \multicolumn{2}{p{0.3\linewidth}}{
        \cellcolor[rgb]{0.9,0.9,0.9}{
            \makecell[{{p{\linewidth}}}]{
                \texttt{\tiny{[GM$|$GM]}}
                \texttt{yes} \\
            }
        }
    }
    & & \\ \\

    \theutterance \stepcounter{utterance}  
    & & & \multicolumn{2}{p{0.3\linewidth}}{
        \cellcolor[rgb]{0.9,0.9,0.9}{
            \makecell[{{p{\linewidth}}}]{
                \texttt{\tiny{[GM$|$GM]}}
                \texttt{game\_result = LOSE} \\
            }
        }
    }
    & & \\ \\

\end{supertabular}
}

\end{document}
